\section{Solvers}

Three forward models coexist, used at different costs for different
questions:

\begin{description}
  \item[Ray model] (\cmd{sim/model/}, used by the studio). Pure
    numpy, interactive speed. Propagates diametral transits through
    the per-grain slowness field to predict travel times and the
    azimuthal velocity pattern. This is the model the inversion
    inverts, and the GUI's default backend.
  \item[Anisotropic FDTD] (\cmd{sim/core/fdtd.py}). The full-wave
    workhorse: a label-indexed elastic finite-difference solver in
    which every cell's stiffness comes from its grain's orientation.
    Runs on GPU through cupy; produces the synthetic pulse-echo
    traces the analysis layer consumes. A pseudospectral solver
    provides an independent numerical scheme for validation.
  \item[Finite-element arbiter] (\cmd{sim/fe\_crosscheck/}). A
    completely independent solver stack (its own mesher, elements
    and probes) used to adjudicate the FDTD results. When two
    unrelated discretisations of the same physics agree, the result
    is a property of the physics rather than the scheme.
\end{description}

The rotation machinery rotates the specimen rather than the probe,
matching the physical machine, and the bit-equivalence gate
(\cmd{validate\_labels.py}) guarantees the rotated label field feeds
the solver exactly the specimen the builder produced.
